% answers.tex —— 2012 年·广东卷·理科 · 参考答案与详细解析片段
% 由 tools/generate_projects.py 依来源总账自 2012/2012_数学真题_广东理_试题.md + 2012/2012_数学真题_广东理_解析正文.md + 2012/2012_数学真题_广东理_答案要点.md 生成，勿手改（改源 MD 后重跑）


\begin{answer}
\gksolhead{1}{}

\textbf{题目}：

设 $\mathrm i$ 为虚数单位，则复数 $\dfrac{5-6\mathrm i}{\mathrm i}=$\mbox{（　　）}

（A）$6+5\mathrm i$；（B）$6-5\mathrm i$；（C）$-6+5\mathrm i$；（D）$-6-5\mathrm i$。

\Needspace{4\baselineskip}
\textbf{解析}

\Needspace{5\baselineskip}
由 $\dfrac1{\mathrm i}=\dfrac{\mathrm i}{\mathrm i^2}=-\mathrm i$，把除法化为乘法：
\nopagebreak[4]
\[\frac{5-6\mathrm i}{\mathrm i}=(5-6\mathrm i)\cdot(-\mathrm i)=-5\mathrm i+6\mathrm i^2.\]
\Needspace{5\baselineskip}
代入 $\mathrm i^2=-1$，得
\nopagebreak[4]
\[-5\mathrm i-6=-6-5\mathrm i.\]

选 D。

\gkanswer{答案：D。$\dfrac1{\mathrm i}=-\mathrm i$，故 $\dfrac{5-6\mathrm i}{\mathrm i}=(5-6\mathrm i)(-\mathrm i)=-6-5\mathrm i$。}

\end{answer}
\begin{answer}
\gksolhead{2}{}

\textbf{题目}：

设集合 $U=\{1,2,3,4,5,6\}$，$M=\{1,2,4\}$，则 $\complement_U M=$\mbox{（　　）}

（A）$U$；（B）$\{1,3,5\}$；（C）$\{3,5,6\}$；（D）$\{2,4,6\}$。

\Needspace{4\baselineskip}
\textbf{解析}

\Needspace{5\baselineskip}
$\complement_U M$ 由 $U$ 中不属于 $M$ 的元素组成。逐个核对 $U=\{1,2,3,4,5,6\}$ 与 $M=\{1,2,4\}$：
\nopagebreak[4]
\[1,2,4\in M;\qquad 3,5,6\notin M.\]
\Needspace{5\baselineskip}
故
\nopagebreak[4]
\[\complement_U M=\{3,5,6\}.\]

选 C。

\gkanswer{答案：C。$U$ 中去掉 $M=\{1,2,4\}$ 的元素，剩 $\{3,5,6\}$。}

\end{answer}
\begin{answer}
\gksolhead{3}{}

\textbf{题目}：

若向量 $\overrightarrow{BA}=(2,3)$，$\overrightarrow{CA}=(4,7)$，则 $\overrightarrow{BC}=$\mbox{（　　）}

（A）$(-2,-4)$；（B）$(3,4)$；（C）$(6,10)$；（D）$(-6,-10)$。

\Needspace{4\baselineskip}
\textbf{解析}

\Needspace{5\baselineskip}
由向量加法的三角形法则，
\nopagebreak[4]
\[\overrightarrow{BC}=\overrightarrow{BA}+\overrightarrow{AC}.\]
\Needspace{5\baselineskip}
又 $\overrightarrow{AC}=-\overrightarrow{CA}$，故
\nopagebreak[4]
\[\overrightarrow{BC}=\overrightarrow{BA}-\overrightarrow{CA}=(2,3)-(4,7)=(2-4,\,3-7)=(-2,-4).\]

选 A。

\gkanswer{答案：A。$\overrightarrow{BC}=\overrightarrow{BA}+\overrightarrow{AC}=\overrightarrow{BA}-\overrightarrow{CA}=(2-4,3-7)=(-2,-4)$。}

\end{answer}
\begin{answer}
\gksolhead{4}{}

\textbf{题目}：

下列函数中，在区间 $(0,+\infty)$ 上为增函数的是\mbox{（　　）}

（A）$y=\ln(x+2)$；（B）$y=-\sqrt{x+1}$；（C）$y=\left(\dfrac12\right)^x$；（D）$y=x+\dfrac1x$。

\Needspace{4\baselineskip}
\textbf{解析}

逐项判断：

（A）$x+2$ 在 $(0,+\infty)$ 上为正且递增，$\ln u$ 在 $u>0$ 上递增，复合后 $y=\ln(x+2)$ 递增；

（B）$\sqrt{x+1}$ 递增，添负号后 $y=-\sqrt{x+1}$ 递减；

（C）底数 $\dfrac12\in(0,1)$，$y=\left(\dfrac12\right)^x$ 递减；

（D）$y'=1-\dfrac1{x^2}$，在 $(0,1)$ 上为负、在 $(1,+\infty)$ 上为正，先减后增，不是单调函数。

只有（A）符合。

选 A。

\gkanswer{答案：A。$\ln(x+2)$ 在 $(0,+\infty)$ 上递增；其余三式分别递减或先减后增。}

\end{answer}
\begin{answer}
\gksolhead{5}{}

\textbf{题目}：

\Needspace{5\baselineskip}
已知变量 $x$，$y$ 满足约束条件
\nopagebreak[4]
\[\begin{cases}
y\leqslant 2,\\
x+y\geqslant 1,\\
x-y\leqslant 1,
\end{cases}\]
则 $z=3x+y$ 的最大值为\mbox{（　　）}

（A）$12$；（B）$11$；（C）$3$；（D）$-1$。

\Needspace{4\baselineskip}
\textbf{解析}

\Needspace{5\baselineskip}
由 $x-y\leqslant1$ 得 $x\leqslant y+1$。于是
\nopagebreak[4]
\[z=3x+y\leqslant3(y+1)+y=4y+3.\]
\Needspace{5\baselineskip}
又 $y\leqslant2$，故
\nopagebreak[4]
\[z\leqslant4\times2+3=11.\]
\Needspace{5\baselineskip}
当 $x=3$，$y=2$ 时，$x+y=5\geqslant1$，$x-y=1\leqslant1$，$y=2$，三条件均满足，且
\nopagebreak[4]
\[z=3\times3+2=11.\]
所以最大值为 $11$。

选 B。

\gkanswer{答案：B。由 $x\leqslant y+1$、$y\leqslant2$ 得 $z=3x+y\leqslant4y+3\leqslant11$，在 $(3,2)$ 处取到。}

\end{answer}
\begin{answer}
\gksolhead{6}{}

\textbf{题目}：

某几何体的三视图如图所示，则它的体积为\mbox{（　　）}
\nopagebreak[4]
\begin{center}% ============================================================
% f2012_guangdong_science_q06_views.tex
% 2012 广东卷（理科）第 6 题 三视图（TikZ 矢量重绘）
% 正视图、侧视图：下方矩形 6×5，上方等腰三角形底 6、腰 5（高 4）
% 俯视图：直径 6 的圆；还原为圆柱（r=3, h=5）与圆锥（r=3, 母线 5, h=4）
% ============================================================
\begin{tikzpicture}[
  x=1cm, y=1cm, line width=0.6pt, font=\footnotesize, >=Stealth,
  dim/.style={<->, line width=0.4pt}
]

  % ---------------- 正视图 ----------------
  \begin{scope}
    \draw (-1.5,0) rectangle (1.5,2.5);
    \draw (-1.5,2.5) -- (0,4.5) -- (1.5,2.5);
    % 底边 6
    \draw[dim] (-1.5,-0.42) -- (1.5,-0.42);
    \draw (-1.5,-0.34) -- (-1.5,-0.50);
    \draw (1.5,-0.34) -- (1.5,-0.50);
    \node[below] at (0,-0.42) {$6$};
    % 高 5
    \draw[dim] (1.88,0) -- (1.88,2.5);
    \draw (1.80,0) -- (1.96,0);
    \draw (1.80,2.5) -- (1.96,2.5);
    \node[right] at (1.90,1.25) {$5$};
    % 两腰 5
    \node[rotate=53.13] at (-1.20,3.30) {$5$};
    \node[rotate=-53.13] at (1.20,3.30) {$5$};
    \node at (0,-1.10) {正视图};
  \end{scope}

  % ---------------- 侧视图 ----------------
  \begin{scope}[xshift=5.4cm]
    \draw (-1.5,0) rectangle (1.5,2.5);
    \draw (-1.5,2.5) -- (0,4.5) -- (1.5,2.5);
    \draw[dim] (-1.5,-0.42) -- (1.5,-0.42);
    \draw (-1.5,-0.34) -- (-1.5,-0.50);
    \draw (1.5,-0.34) -- (1.5,-0.50);
    \node[below] at (0,-0.42) {$6$};
    \draw[dim] (1.88,0) -- (1.88,2.5);
    \draw (1.80,0) -- (1.96,0);
    \draw (1.80,2.5) -- (1.96,2.5);
    \node[right] at (1.90,1.25) {$5$};
    \node[rotate=53.13] at (-1.20,3.30) {$5$};
    \node[rotate=-53.13] at (1.20,3.30) {$5$};
    \node at (0,-1.10) {侧视图};
  \end{scope}

  % ---------------- 俯视图 ----------------
  \begin{scope}[yshift=-4.5cm]
    \draw (0,0) circle (1.5);
    \fill (0,0) circle (1.1pt);
    \node at (0,-1.95) {俯视图};
  \end{scope}

\end{tikzpicture}
\end{center}
\nopagebreak[4]
（A）$12\pi$；（B）$45\pi$；（C）$57\pi$；（D）$81\pi$。

\Needspace{4\baselineskip}
\textbf{解析}

由三视图还原几何体：

\Needspace{5\baselineskip}
正视图与侧视图的下部都是宽 $6$、高 $5$ 的矩形，上部都是底边 $6$、两腰 $5$ 的等腰三角形；俯视图是直径 $6$ 的圆。因此该几何体是底面半径
\nopagebreak[4]
\[r=\frac62=3\]
的圆柱，上面接一个同底圆锥。

\Needspace{5\baselineskip}
圆锥的母线为 $5$，底面半径为 $3$，故圆锥的高为
\nopagebreak[4]
\[h_1=\sqrt{5^2-3^2}=\sqrt{16}=4.\]
\Needspace{5\baselineskip}
圆柱的高为 $h_2=5$。体积为
\nopagebreak[4]
\[\begin{aligned}
V&=V_{\text{\text{圆柱}}}+V_{\text{\text{圆锥}}}\\
&=\pi r^2h_2+\frac13\pi r^2h_1\\
&=\pi\cdot3^2\cdot5+\frac13\pi\cdot3^2\cdot4\\
&=45\pi+12\pi=57\pi.
\end{aligned}\]

选 C。

\gkanswer{答案：C。几何体为圆柱与圆锥的组合，$r=3$，圆柱高 $5$，圆锥高 $\sqrt{5^2-3^2}=4$。}

\end{answer}
\begin{answer}
\gksolhead{7}{}

\textbf{题目}：

从个位数与十位数之和为奇数的两位数中任取一个，其个位数为 $0$ 的概率是\mbox{（　　）}

（A）$\dfrac49$；（B）$\dfrac13$；（C）$\dfrac29$；（D）$\dfrac19$。

\Needspace{4\baselineskip}
\textbf{解析}

设两位数的十位数字为 $a\in\{1,2,\ldots,9\}$，个位数字为 $b\in\{0,1,\ldots,9\}$。条件 $a+b$ 为奇数，即 $a$、$b$ 一奇一偶。

十位为奇数（$5$ 种）、个位为偶数（$5$ 种）：$5\times5=25$ 个；

十位为偶数（$4$ 种）、个位为奇数（$5$ 种）：$4\times5=20$ 个。

\Needspace{5\baselineskip}
合计
\nopagebreak[4]
\[25+20=45\]
个两位数。

\Needspace{5\baselineskip}
其中个位数为 $0$ 的，必须是十位为奇数、个位为偶数且 $b=0$，即
\nopagebreak[4]
\[10,30,50,70,90,\]
\Needspace{5\baselineskip}
共 $5$ 个。所求概率为
\nopagebreak[4]
\[P=\frac5{45}=\frac19.\]

选 D。

\gkanswer{答案：D。两位数字之和为奇数者共 $45$ 个，其中个位为 $0$ 的有 $5$ 个，概率 $\dfrac5{45}=\dfrac19$。}

\end{answer}
\begin{answer}
\gksolhead{8}{}

\textbf{题目}：

对任意两个非零的平面向量 $\boldsymbol\alpha$ 和 $\boldsymbol\beta$，定义 $\boldsymbol\alpha\circ\boldsymbol\beta=\dfrac{\boldsymbol\alpha\cdot\boldsymbol\beta}{\boldsymbol\beta\cdot\boldsymbol\beta}$。若平面向量 $\boldsymbol a$，$\boldsymbol b$ 满足 $|\boldsymbol a|\geqslant|\boldsymbol b|>0$，$\boldsymbol a$ 与 $\boldsymbol b$ 的夹角 $\theta\in\left(0,\dfrac\pi4\right)$，且 $\boldsymbol a\circ\boldsymbol b$ 和 $\boldsymbol b\circ\boldsymbol a$ 都在集合 $\left\{\dfrac n2\mid n\in\mathbf Z\right\}$ 中，则 $\boldsymbol a\circ\boldsymbol b=$\mbox{（　　）}

（A）$\dfrac12$；（B）$1$；（C）$\dfrac32$；（D）$\dfrac52$。

\Needspace{4\baselineskip}
\textbf{解析}

\Needspace{5\baselineskip}
记 $|\boldsymbol a|=a$，$|\boldsymbol b|=b$，夹角为 $\theta$，则
\nopagebreak[4]
\[\boldsymbol a\circ\boldsymbol b=\frac{\boldsymbol a\cdot\boldsymbol b}{\boldsymbol b\cdot\boldsymbol b}
=\frac{ab\cos\theta}{b^2}=\frac ab\cos\theta,\]
\nopagebreak[4]
\[\boldsymbol b\circ\boldsymbol a=\frac{\boldsymbol a\cdot\boldsymbol b}{\boldsymbol a\cdot\boldsymbol a}
=\frac{ab\cos\theta}{a^2}=\frac ba\cos\theta.\]

\Needspace{5\baselineskip}
令 $r=\dfrac ab$。由 $a\geqslant b>0$ 得 $r\geqslant1$；记 $c=\cos\theta$。因 $\theta\in\left(0,\dfrac\pi4\right)$，
\nopagebreak[4]
\[c\in\left(\frac{\sqrt2}2,1\right).\]
\Needspace{5\baselineskip}
于是
\nopagebreak[4]
\[\boldsymbol b\circ\boldsymbol a=\frac cr,\qquad
\boldsymbol a\circ\boldsymbol b=rc.\]

\Needspace{5\baselineskip}
先看 $\boldsymbol b\circ\boldsymbol a$。由 $r\geqslant1$、$0<c<1$ 得
\nopagebreak[4]
\[0<\frac cr<1.\]
\Needspace{5\baselineskip}
它属于 $\left\{\dfrac n2\mid n\in\mathbf Z\right\}$，即必须是正的半整数；小于 $1$ 的正半整数只有 $\dfrac12$，故
\nopagebreak[4]
\[\frac cr=\frac12,\qquad r=2c.\]

\Needspace{5\baselineskip}
再算
\nopagebreak[4]
\[\boldsymbol a\circ\boldsymbol b=rc=2c^2.\]
\Needspace{5\baselineskip}
由 $c^2\in\left(\dfrac12,1\right)$ 得
\nopagebreak[4]
\[2c^2\in(1,2).\]
\Needspace{5\baselineskip}
它是半整数，开区间 $(1,2)$ 内唯一的半整数是 $\dfrac32$，故
\nopagebreak[4]
\[\boldsymbol a\circ\boldsymbol b=\frac32.\]
此时 $c^2=\dfrac34$，$c=\dfrac{\sqrt3}2$，$r=\sqrt3\geqslant1$，条件可以同时成立。

选 C。

\gkanswer{答案：C。记 $r=\dfrac{|\boldsymbol a|}{|\boldsymbol b|}\geqslant1$，$c=\cos\theta\in\left(\dfrac{\sqrt2}2,1\right)$，则 $\boldsymbol b\circ\boldsymbol a=\dfrac cr=\dfrac12$，$\boldsymbol a\circ\boldsymbol b=rc=2c^2=\dfrac32$。}

\end{answer}
\begin{answer}
\gksolhead{9}{}

\textbf{题目}：

不等式 $|x+2|-|x|\leqslant 1$ 的解集为\gkblank{}。

\Needspace{4\baselineskip}
\textbf{解析}

按零点 $-2$、$0$ 分段讨论。

\Needspace{5\baselineskip}
当 $x<-2$ 时，$x+2<0$，$x<0$，故
\nopagebreak[4]
\[|x+2|-|x|=-(x+2)-(-x)=-2\leqslant1,\]
恒成立。

\Needspace{5\baselineskip}
当 $-2\leqslant x<0$ 时，$x+2\geqslant0$，$x<0$，故
\nopagebreak[4]
\[|x+2|-|x|=(x+2)-(-x)=2x+2.\]
\Needspace{5\baselineskip}
不等式化为
\nopagebreak[4]
\[2x+2\leqslant1,\qquad x\leqslant-\frac12.\]
\Needspace{5\baselineskip}
与 $-2\leqslant x<0$ 取交集，得
\nopagebreak[4]
\[-2\leqslant x\leqslant-\frac12.\]

\Needspace{5\baselineskip}
当 $x\geqslant0$ 时，$x+2>0$，$x\geqslant0$，故
\nopagebreak[4]
\[|x+2|-|x|=(x+2)-x=2,\]
而 $2\leqslant1$ 不成立，此段无解。

\Needspace{5\baselineskip}
合并前两段：
\nopagebreak[4]
\[\left(-\infty,-\frac12\right].\]
代回检验：$x=-\dfrac12$ 时 $\left|\dfrac32\right|-\left|-\dfrac12\right|=1$，取等号；$x=-3$ 时 $1-3=-2\leqslant1$，均符合。

\gkanswer{答案：$\left(-\infty,-\dfrac12\right]$。分 $x<-2$、$-2\leqslant x<0$、$x\geqslant0$ 三段讨论，仅前两段的部分满足。}

\end{answer}
\begin{answer}
\gksolhead{10}{}

\textbf{题目}：

$\left(x^2+\dfrac1x\right)^6$ 的展开式中 $x^3$ 的系数为\gkblank{}。（用数字作答）

\Needspace{4\baselineskip}
\textbf{解析}

\Needspace{5\baselineskip}
二项展开式的通项为
\nopagebreak[4]
\[T_{k+1}=\mathrm C_6^k\left(x^2\right)^{6-k}\left(\frac1x\right)^k
=\mathrm C_6^k x^{2(6-k)-k}
=\mathrm C_6^k x^{12-3k},\qquad k=0,1,\ldots,6.\]
\Needspace{5\baselineskip}
要求 $x^3$ 项，令
\nopagebreak[4]
\[12-3k=3,\]
\Needspace{5\baselineskip}
解得
\nopagebreak[4]
\[k=3.\]
\Needspace{5\baselineskip}
故该项系数为
\nopagebreak[4]
\[\mathrm C_6^3=\frac{6\times5\times4}{3\times2\times1}=20.\]

\gkanswer{答案：$20$。通项为 $\mathrm C_6^k x^{12-3k}$，令 $12-3k=3$ 得 $k=3$，系数 $\mathrm C_6^3=20$。}

\end{answer}
\begin{answer}
\gksolhead{11}{}

\textbf{题目}：

已知递增的等差数列 $\{a_n\}$ 满足 $a_1=1$，$a_3=a_2^2-4$，则 $a_n=$\gkblank{}。

\Needspace{4\baselineskip}
\textbf{解析}

\Needspace{5\baselineskip}
设等差数列的公差为 $d$。因数列递增，$d>0$。由 $a_1=1$ 得
\nopagebreak[4]
\[a_2=1+d,\qquad a_3=1+2d.\]
\Needspace{5\baselineskip}
代入 $a_3=a_2^2-4$：
\nopagebreak[4]
\[1+2d=(1+d)^2-4=d^2+2d-3.\]
\Needspace{5\baselineskip}
两边消去 $2d$，得
\nopagebreak[4]
\[1=d^2-3,\qquad d^2=4.\]
\Needspace{5\baselineskip}
因 $d>0$，取
\nopagebreak[4]
\[d=2.\]
\Needspace{5\baselineskip}
于是
\nopagebreak[4]
\[a_n=1+2(n-1)=2n-1.\]
检验：$a_2=3$，$a_3=5$，而 $3^2-4=5$，符合。

\gkanswer{答案：$a_n=2n-1$。设公差 $d>0$，由 $1+2d=(1+d)^2-4$ 得 $d=2$。}

\end{answer}
\begin{answer}
\gksolhead{12}{}

\textbf{题目}：

曲线 $y=x^3-x+3$ 在点 $(1,3)$ 处的切线方程为\gkblank{}。

\Needspace{4\baselineskip}
\textbf{解析}

\Needspace{5\baselineskip}
设 $y=f(x)=x^3-x+3$，则
\nopagebreak[4]
\[f'(x)=3x^2-1.\]
\Needspace{5\baselineskip}
在 $x=1$ 处，切线的斜率
\nopagebreak[4]
\[k=f'(1)=3-1=2.\]
\Needspace{5\baselineskip}
切点为 $(1,3)$，由点斜式
\nopagebreak[4]
\[y-3=2(x-1),\]
\Needspace{5\baselineskip}
即
\nopagebreak[4]
\[y=2x+1.\]

\gkanswer{答案：$y=2x+1$。$y'=3x^2-1$，在 $x=1$ 处斜率为 $2$。}

\end{answer}
\begin{answer}
\gksolhead{13}{}

\textbf{题目}：

执行如图所示的程序框图，若输入 $n$ 的值为 $8$，则输出 $s$ 的值为\gkblank{}。
\nopagebreak[4]
\begin{center}% ============================================================
% f2012_guangdong_science_q13_flowchart.tex
% 2012 广东卷（理科）第 13 题 程序框图（TikZ 矢量重绘）
% 开始 → 输入 n → i=2, k=1, s=1 → 判断 i<n
%   是：s=(1/k)(s×i) → i=i+2 → k=k+1 → 回到判断
%   否：输出 s → 结束
% ============================================================
\begin{tikzpicture}[
  x=1cm, y=1cm, line width=0.6pt, font=\footnotesize, >=Stealth,
  every node/.style={inner sep=3pt},
  proc/.style={rectangle, draw, minimum width=2.9cm, minimum height=0.68cm},
  io/.style={trapezium, draw, trapezium left angle=70, trapezium right angle=110,
             minimum width=3.0cm, minimum height=0.72cm, trapezium stretches=true},
  dec/.style={diamond, draw, aspect=2.3, minimum width=2.5cm, minimum height=0.95cm},
  term/.style={rectangle, rounded corners=8pt, draw, minimum width=1.9cm, minimum height=0.60cm}
]

  \node[term] (start) at (0,0) {开始};
  \node[io]   (input) at (0,-1.25) {输入 $n$};
  \node[proc] (init)  at (0,-2.50) {$i=2$，$k=1$，$s=1$};
  \node[dec]  (judge) at (0,-3.95) {$i<n$};

  \node[proc, minimum height=1.10cm] (body) at (0,-5.55) {$s=\dfrac1k(s\times i)$};
  \node[proc] (inc)   at (0,-6.95) {$i=i+2$};
  \node[proc] (kk)    at (0,-8.20) {$k=k+1$};

  \node[io]   (out)   at (3.90,-5.55) {输出 $s$};
  \node[term] (end)   at (3.90,-6.95) {结束};

  \draw[->] (start) -- (input);
  \draw[->] (input) -- (init);
  \draw[->] (init) -- (judge);
  \draw[->] (judge) -- node[right] {是} (body);
  \draw[->] (body) -- (inc);
  \draw[->] (inc) -- (kk);
  \draw[->] (kk.west) -- (-2.55,-8.20) -- (-2.55,-3.95) -- (judge.west);
  \draw[->] (judge.east) -- (3.90,-3.95) -- (out.north);
  \node at (2.55,-3.72) {否};
  \draw[->] (out) -- (end);

\end{tikzpicture}
\end{center}
\nopagebreak[4]
\Needspace{4\baselineskip}
\textbf{解析}

\Needspace{5\baselineskip}
按框图执行，初值为
\nopagebreak[4]
\[i=2,\qquad k=1,\qquad s=1.\]
输入 $n=8$。

\Needspace{5\baselineskip}
第一次判断：$i=2<8$ 成立，
\nopagebreak[4]
\[s=\frac1k(s\times i)=\frac11(1\times2)=2,\qquad i=2+2=4,\qquad k=1+1=2.\]

\Needspace{5\baselineskip}
第二次判断：$i=4<8$ 成立，
\nopagebreak[4]
\[s=\frac12(2\times4)=4,\qquad i=4+2=6,\qquad k=2+1=3.\]

\Needspace{5\baselineskip}
第三次判断：$i=6<8$ 成立，
\nopagebreak[4]
\[s=\frac13(4\times6)=8,\qquad i=6+2=8,\qquad k=3+1=4.\]

\Needspace{5\baselineskip}
第四次判断：$i=8<8$ 不成立，退出循环，输出
\nopagebreak[4]
\[s=8.\]

\gkanswer{答案：$8$。依次得 $s=2,4,8$，对应 $i=2,4,6$；$i=8$ 时退出循环。}

\end{answer}
\begin{answer}
\gksolhead{14}{}

\textbf{题目}：

\Needspace{5\baselineskip}
在平面直角坐标系 $xOy$ 中，曲线 $C_1$ 和 $C_2$ 的参数方程分别为
\nopagebreak[4]
\[\begin{cases}
x=t,\\
y=\sqrt t,
\end{cases}\quad(t\text{ \text{为参数}})\]
\Needspace{5\baselineskip}
和
\nopagebreak[4]
\[\begin{cases}
x=\sqrt2\cos\theta,\\
y=\sqrt2\sin\theta,
\end{cases}\quad(\theta\text{ \text{为参数}}),\]
则曲线 $C_1$ 与 $C_2$ 的交点坐标为\gkblank{}。

\Needspace{4\baselineskip}
\textbf{解析}

\Needspace{5\baselineskip}
对 $C_1$，由 $x=t$ 得 $t=x$，代入 $y=\sqrt t$：
\nopagebreak[4]
\[y=\sqrt x,\qquad x\geqslant0,\ y\geqslant0,\]
\Needspace{5\baselineskip}
即
\nopagebreak[4]
\[y^2=x\quad(y\geqslant0).\]

\Needspace{5\baselineskip}
对 $C_2$，两式平方相加：
\nopagebreak[4]
\[x^2+y^2=2\cos^2\theta+2\sin^2\theta=2,\]
故 $C_2$ 是圆 $x^2+y^2=2$。

\Needspace{5\baselineskip}
联立
\nopagebreak[4]
\[\begin{cases}
y^2=x,\\
x^2+y^2=2,
\end{cases}\]
\Needspace{5\baselineskip}
把 $y^2=x$ 代入第二式：
\nopagebreak[4]
\[x^2+x-2=0,\qquad (x+2)(x-1)=0.\]
得 $x=1$ 或 $x=-2$。因 $C_1$ 要求 $x\geqslant0$，舍去 $x=-2$；当 $x=1$ 时 $y^2=1$，又 $y\geqslant0$，故 $y=1$。

\Needspace{5\baselineskip}
交点坐标为
\nopagebreak[4]
\[(1,1).\]
代回检验：$1^2+1^2=2$，且 $1=\sqrt1$，两曲线方程均成立。

\gkanswer{答案：$(1,1)$。$C_1$ 为 $y^2=x$（$y\geqslant0$），$C_2$ 为 $x^2+y^2=2$，联立得 $x=1$。}

\end{answer}
\begin{answer}
\gksolhead{15}{}

\textbf{题目}：

如图，圆 $O$ 的半径为 $1$，$A$，$B$，$C$ 是圆周上的三点，满足 $\angle ABC=30^\circ$，过点 $A$ 作圆 $O$ 的切线与 $OC$ 的延长线交于点 $P$，则 $PA=$\gkblank{}。
\nopagebreak[4]
\begin{center}% ============================================================
% f2012_guangdong_science_q15_circle.tex
% 2012 广东卷（理科）第 15 题 圆与切线（TikZ 矢量重绘）
% 圆 O 半径 1；A、B、C 在圆上，∠ABC=30°；
% 过 A 的切线与 OC 的延长线交于 P。图中 ∠AOC=60°，PA=√3。
% ============================================================
\begin{tikzpicture}[
  x=1cm, y=1cm, line width=0.6pt, font=\footnotesize, >=Stealth
]
  \coordinate (O) at (0,0);
  \coordinate (A) at (60:1.5);
  \coordinate (C) at (0:1.5);
  \coordinate (B) at (220:1.5);
  \coordinate (P) at (3.0,0);

  \draw (O) circle (1.5);
  \draw (A) -- (B) -- (C);
  \draw (O) -- (A);
  \draw (O) -- (P);
  \draw (A) -- (P);

  \fill (O) circle (1.0pt);
  \node[above left] at (O) {$O$};
  \node[above right] at (A) {$A$};
  \node[below left] at (B) {$B$};
  \node[below] at (C) {$C$};
  \node[right] at (P) {$P$};
\end{tikzpicture}
\end{center}
\nopagebreak[4]
\Needspace{4\baselineskip}
\textbf{解析}

\Needspace{5\baselineskip}
$\angle ABC$ 是圆周角，它所对的弧为弧 $AC$（不含 $B$）；同弧所对的圆心角为 $\angle AOC$，故
\nopagebreak[4]
\[\angle AOC=2\angle ABC=2\times30^\circ=60^\circ.\]

\Needspace{5\baselineskip}
$PA$ 与圆相切于 $A$，故
\nopagebreak[4]
\[OA\perp PA.\]
\Needspace{5\baselineskip}
点 $P$ 在 $OC$ 的延长线上，所以射线 $OP$ 与射线 $OC$ 重合，
\nopagebreak[4]
\[\angle AOP=\angle AOC=60^\circ.\]

\Needspace{5\baselineskip}
在直角三角形 $OAP$ 中，$\angle OAP=90^\circ$，$OA=1$，于是
\nopagebreak[4]
\[\tan\angle AOP=\frac{PA}{OA},\]
\Needspace{5\baselineskip}
即
\nopagebreak[4]
\[PA=OA\tan60^\circ=1\times\sqrt3=\sqrt3.\]

\gkanswer{答案：$\sqrt3$。$\angle AOC=2\angle ABC=60^\circ$，在直角三角形 $OAP$ 中 $PA=OA\tan60^\circ=\sqrt3$。}

\end{answer}
\begin{answer}
\gksolhead{16}{}

\textbf{题目}：

已知函数 $f(x)=2\cos\left(\omega x+\dfrac\pi6\right)$（其中 $\omega>0$，$x\in\mathbf R$）的最小正周期为 $10\pi$。

（1） 求 $\omega$ 的值；

（2） 设 $\alpha$，$\beta\in\left[0,\dfrac\pi2\right]$，$f\left(5\alpha+\dfrac{5\pi}{3}\right)=-\dfrac65$，$f\left(5\beta-\dfrac{5\pi}{6}\right)=\dfrac{16}{17}$，求 $\cos(\alpha+\beta)$ 的值。

\Needspace{4\baselineskip}
\textbf{解析}

\textbf{（1）}

\Needspace{5\baselineskip}
函数 $f(x)=2\cos\left(\omega x+\dfrac\pi6\right)$ 的最小正周期为
\nopagebreak[4]
\[T=\frac{2\pi}{\omega}.\]
\Needspace{5\baselineskip}
已知 $T=10\pi$，且 $\omega>0$，故
\nopagebreak[4]
\[\frac{2\pi}{\omega}=10\pi,\qquad \omega=\frac{2\pi}{10\pi}=\frac15.\]

\textbf{（2）}

\Needspace{5\baselineskip}
先化简两个条件中的角。对第一个条件，
\nopagebreak[4]
\[\frac15\left(5\alpha+\frac{5\pi}{3}\right)+\frac\pi6
=\alpha+\frac\pi3+\frac\pi6
=\alpha+\frac\pi2.\]
\Needspace{5\baselineskip}
故
\nopagebreak[4]
\[f\left(5\alpha+\frac{5\pi}{3}\right)
=2\cos\left(\alpha+\frac\pi2\right)
=-2\sin\alpha=-\frac65,\]
\Needspace{5\baselineskip}
从而
\nopagebreak[4]
\[\sin\alpha=\frac35.\]
\Needspace{5\baselineskip}
因 $\alpha\in\left[0,\dfrac\pi2\right]$，取正值
\nopagebreak[4]
\[\cos\alpha=\sqrt{1-\frac9{25}}=\frac45.\]

\Needspace{5\baselineskip}
对第二个条件，
\nopagebreak[4]
\[\frac15\left(5\beta-\frac{5\pi}{6}\right)+\frac\pi6
=\beta-\frac\pi6+\frac\pi6=\beta.\]
\Needspace{5\baselineskip}
故
\nopagebreak[4]
\[f\left(5\beta-\frac{5\pi}{6}\right)=2\cos\beta=\frac{16}{17},\]
\Needspace{5\baselineskip}
即
\nopagebreak[4]
\[\cos\beta=\frac8{17}.\]
\Needspace{5\baselineskip}
因 $\beta\in\left[0,\dfrac\pi2\right]$，
\nopagebreak[4]
\[\sin\beta=\sqrt{1-\frac{64}{289}}=\sqrt{\frac{225}{289}}=\frac{15}{17}.\]

\Needspace{5\baselineskip}
由两角和的余弦公式，
\nopagebreak[4]
\[\begin{aligned}
\cos(\alpha+\beta)
&=\cos\alpha\cos\beta-\sin\alpha\sin\beta\\
&=\frac45\cdot\frac8{17}-\frac35\cdot\frac{15}{17}\\
&=\frac{32}{85}-\frac{45}{85}\\
&=-\frac{13}{85}.
\end{aligned}\]

\gkanswer{答案（1）：$\omega=\dfrac15$。由 $\dfrac{2\pi}{\omega}=10\pi$ 得。}

\gkanswer{答案（2）：$-\dfrac{13}{85}$。由条件得 $\sin\alpha=\dfrac35$，$\cos\alpha=\dfrac45$，$\cos\beta=\dfrac8{17}$，$\sin\beta=\dfrac{15}{17}$，再用两角和的余弦公式。}

\end{answer}
\begin{answer}
\gksolhead{17}{}

\textbf{题目}：

某班 $50$ 位学生期中考试数学成绩的频率分布直方图如图所示，其中成绩分组区间是：$[40,50)$，$[50,60)$，$[60,70)$，$[70,80)$，$[80,90)$，$[90,100]$。
\nopagebreak[4]
\begin{center}% ============================================================
% f2012_guangdong_science_q17_histogram.tex
% 2012 广东卷（理科）第 17 题 频率分布直方图（TikZ 矢量重绘）
% 组距 10；各组 [40,50)…[90,100] 的“频率/组距”依次为
% 0.006, 0.006, 0.010, 0.054, x, 0.006，其中 x=0.018
% ============================================================
\begin{tikzpicture}[
  x=1cm, y=1cm, line width=0.6pt, font=\footnotesize, >=Stealth
]
  \path[draw=none] (-0.75,-0.75) rectangle (8.05,7.15);
  % 坐标轴
  \draw[->] (-0.35,0) -- (7.55,0) node[below] {成绩};
  \draw[->] (0,-0.05) -- (0,6.55);
  \node[right, inner sep=3pt] at (0.10,6.05) {$\dfrac{\text{频率}}{\text{组距}}$};
  % 断轴符号
  \draw (-0.30,0) -- (-0.20,0.15) -- (-0.10,0) -- (0.00,0.15);

  % 六个矩形，高度按 100 倍纵坐标绘制
  \draw (0.5,0) rectangle (1.5,0.6);
  \draw (1.5,0) rectangle (2.5,0.6);
  \draw (2.5,0) rectangle (3.5,1.0);
  \draw (3.5,0) rectangle (4.5,5.4);
  \draw (4.5,0) rectangle (5.5,1.8);
  \draw (5.5,0) rectangle (6.5,0.6);

  % 虚线参考线
  \draw[dashed, line width=0.4pt] (0,5.4) -- (3.5,5.4);
  \draw[dashed, line width=0.4pt] (0,1.8) -- (4.5,1.8);
  \draw[dashed, line width=0.4pt] (0,1.0) -- (2.5,1.0);
  \draw[dashed, line width=0.4pt] (0,0.6) -- (5.5,0.6);

  % 纵轴刻度值
  \node[left] at (0,5.4) {$0.054$};
  \node[left] at (0,1.8) {$x$};
  \node[left] at (0,1.0) {$0.010$};
  \node[left] at (0,0.6) {$0.006$};

  % 横轴刻度
  \draw (0,0) -- (0,-0.09);
  \node[below] at (0,-0.10) {$0$};
  \foreach \xx/\lab in {0.5/40,1.5/50,2.5/60,3.5/70,4.5/80,5.5/90,6.5/100}{
    \draw (\xx,0) -- (\xx,-0.09);
    \node[below] at (\xx,-0.10) {$\lab$};
  }

\end{tikzpicture}
\end{center}
\nopagebreak[4]
（1） 求图中 $x$ 的值；

（2） 从成绩不低于 $80$ 分的学生中随机选取 $2$ 人，该 $2$ 人中成绩在 $90$ 分以上（含 $90$ 分）的人数记为 $\xi$，求 $\xi$ 的数学期望。

\Needspace{4\baselineskip}
\textbf{解析}

组距为 $10$，各组的频率等于矩形的高乘以组距 $10$，六组频率之和为 $1$。

\textbf{（1）}

\Needspace{5\baselineskip}
六个矩形的高依次为
\nopagebreak[4]
\[0.006,\quad0.006,\quad0.010,\quad0.054,\quad x,\quad0.006.\]
\Needspace{5\baselineskip}
由频率之和为 $1$：
\nopagebreak[4]
\[10(0.006+0.006+0.010+0.054+x+0.006)=1.\]
\Needspace{5\baselineskip}
括号内先合并已知数：
\nopagebreak[4]
\[0.006+0.006+0.010+0.054+0.006=0.082,\]
\Needspace{5\baselineskip}
于是
\nopagebreak[4]
\[10(0.082+x)=1,\qquad 0.082+x=0.1,\]
\Needspace{5\baselineskip}
解得
\nopagebreak[4]
\[x=0.018.\]

\textbf{（2）}

\Needspace{5\baselineskip}
成绩不低于 $80$ 分的学生分布在 $[80,90)$ 与 $[90,100]$ 两组，人数分别为
\nopagebreak[4]
\[50\times0.018\times10=9,\qquad 50\times0.006\times10=3,\]
\Needspace{5\baselineskip}
合计
\nopagebreak[4]
\[9+3=12\]
人。

\Needspace{5\baselineskip}
从这 $12$ 人中随机选 $2$ 人，记其中成绩不低于 $90$ 分的人数为 $\xi$。$12$ 人中有 $3$ 人成绩不低于 $90$ 分，$9$ 人成绩在 $[80,90)$ 内，故 $\xi$ 服从超几何分布：
\nopagebreak[4]
\[P(\xi=0)=\frac{\mathrm C_9^2}{\mathrm C_{12}^2}=\frac{36}{66}=\frac6{11},\]
\nopagebreak[4]
\[P(\xi=1)=\frac{\mathrm C_9^1\mathrm C_3^1}{\mathrm C_{12}^2}=\frac{27}{66}=\frac9{22},\]
\nopagebreak[4]
\[P(\xi=2)=\frac{\mathrm C_3^2}{\mathrm C_{12}^2}=\frac3{66}=\frac1{22}.\]
\Needspace{5\baselineskip}
检验概率和为
\nopagebreak[4]
\[\frac6{11}+\frac9{22}+\frac1{22}=\frac{12+9+1}{22}=1.\]
\Needspace{5\baselineskip}
数学期望
\nopagebreak[4]
\[\begin{aligned}
E\xi&=0\cdot\frac6{11}+1\cdot\frac9{22}+2\cdot\frac1{22}\\
&=\frac{9+2}{22}=\frac{11}{22}=\frac12.
\end{aligned}\]
也可直接按比例计算：$E\xi=2\cdot\dfrac3{12}=\dfrac12$。

\gkanswer{答案（1）：$x=0.018$。各组频率之和为 $1$，即 $10(0.082+x)=1$。}

\gkanswer{答案（2）：$\dfrac12$。不低于 $80$ 分者共 $12$ 人，其中不低于 $90$ 分者 $3$ 人，故 $E\xi=2\cdot\dfrac3{12}=\dfrac12$。}

\end{answer}
\begin{answer}
\gksolhead{18}{}

\textbf{题目}：

如图所示，在四棱锥 $P-ABCD$ 中，底面 $ABCD$ 为矩形，$PA\perp$ 平面 $ABCD$，点 $E$ 在线段 $PC$ 上，$PC\perp$ 平面 $BDE$。

（1） 证明：$BD\perp$ 平面 $PAC$；

（2） 若 $PA=1$，$AD=2$，求二面角 $B-PC-A$ 的正切值。
\nopagebreak[4]
\begin{center}% ============================================================
% f2012_guangdong_science_q18_pyramid.tex
% 2012 广东卷（理科）第 18 题 四棱锥 P-ABCD（TikZ 矢量重绘）
% 底面 ABCD 为矩形，PA⊥平面 ABCD，E 在线段 PC 上，PC⊥平面 BDE
% 虚线为不可见棱：PA、AB、AD、AC、BD
% ============================================================
\begin{tikzpicture}[
  x=1cm, y=1cm, line width=0.6pt, font=\footnotesize, >=Stealth
]
  \coordinate (A) at (0,1.1);
  \coordinate (B) at (-2.5,0.1);
  \coordinate (C) at (2.1,0.1);
  \coordinate (D) at (4.6,1.1);
  \coordinate (P) at (-0.7,4.6);
  \coordinate (E) at ($(P)!0.5556!(C)$);

  % 不可见棱
  \draw[dashed, line width=0.45pt] (P) -- (A);
  \draw[dashed, line width=0.45pt] (A) -- (B);
  \draw[dashed, line width=0.45pt] (A) -- (D);
  \draw[dashed, line width=0.45pt] (A) -- (C);
  \draw[dashed, line width=0.45pt] (B) -- (D);

  % 可见棱
  \draw (P) -- (B);
  \draw (P) -- (C);
  \draw (P) -- (D);
  \draw (B) -- (C);
  \draw (C) -- (D);
  \draw (B) -- (E);
  \draw (E) -- (D);

  % 顶点与交点
  \fill (A) circle (1.0pt);
  \fill (B) circle (1.0pt);
  \fill (C) circle (1.0pt);
  \fill (D) circle (1.0pt);
  \fill (P) circle (1.0pt);
  \fill (E) circle (1.0pt);

  \node[above] at (P) {$P$};
  \node[below left] at (A) {$A$};
  \node[below left] at (B) {$B$};
  \node[below] at (C) {$C$};
  \node[right] at (D) {$D$};
  \node[above right] at (E) {$E$};

\end{tikzpicture}
\end{center}
\nopagebreak[4]
\Needspace{4\baselineskip}
\textbf{解析}

\textbf{（1）}

\Needspace{5\baselineskip}
因为 $PA\perp$ 平面 $ABCD$，而 $BD\subset$ 平面 $ABCD$，所以
\nopagebreak[4]
\[PA\perp BD.\]
\Needspace{5\baselineskip}
又 $PC\perp$ 平面 $BDE$，而 $BD\subset$ 平面 $BDE$，所以
\nopagebreak[4]
\[PC\perp BD.\]
\Needspace{5\baselineskip}
$PA$ 与 $PC$ 是平面 $PAC$ 内两条相交直线，且 $PA\cap PC=P$，故
\nopagebreak[4]
\[BD\perp\text{\text{平面} }PAC.\]

\textbf{（2）}

\Needspace{5\baselineskip}
建立空间直角坐标系。以 $A$ 为原点，$\overrightarrow{AB}$ 为 $x$ 轴正方向，$\overrightarrow{AD}$ 为 $y$ 轴正方向，$\overrightarrow{AP}$ 为 $z$ 轴正方向。设 $AB=b$，则
\nopagebreak[4]
\[A(0,0,0),\quad B(b,0,0),\quad D(0,2,0),\quad C(b,2,0),\quad P(0,0,1).\]

\Needspace{5\baselineskip}
由（1）的推导，$PC\perp BD$。计算
\nopagebreak[4]
\[\overrightarrow{PC}=(b,2,-1),\qquad \overrightarrow{BD}=(-b,2,0).\]
\Needspace{5\baselineskip}
垂直条件给出
\nopagebreak[4]
\[\overrightarrow{PC}\cdot\overrightarrow{BD}=-b^2+4=0,\]
\Needspace{5\baselineskip}
故
\nopagebreak[4]
\[b=2.\]
\Needspace{5\baselineskip}
于是
\nopagebreak[4]
\[B(2,0,0),\qquad C(2,2,0).\]

\Needspace{5\baselineskip}
平面 $BPC$ 内有两个向量
\nopagebreak[4]
\[\overrightarrow{PB}=(2,0,-1),\qquad \overrightarrow{PC}=(2,2,-1).\]
\Needspace{5\baselineskip}
取法向量
\nopagebreak[4]
\[\boldsymbol n_1=\overrightarrow{PB}\times\overrightarrow{PC}
=\begin{vmatrix}
\mathbf i&\mathbf j&\mathbf k\\
2&0&-1\\
2&2&-1
\end{vmatrix}
=(2,0,4),\]
\Needspace{5\baselineskip}
可简取
\nopagebreak[4]
\[\boldsymbol n_1=(1,0,2).\]

\Needspace{5\baselineskip}
平面 $APC$ 内有两个向量
\nopagebreak[4]
\[\overrightarrow{AP}=(0,0,1),\qquad \overrightarrow{AC}=(2,2,0).\]
\Needspace{5\baselineskip}
取法向量
\nopagebreak[4]
\[\boldsymbol n_2=\overrightarrow{AP}\times\overrightarrow{AC}
=\begin{vmatrix}
\mathbf i&\mathbf j&\mathbf k\\
0&0&1\\
2&2&0
\end{vmatrix}
=(-2,2,0),\]
\Needspace{5\baselineskip}
可简取
\nopagebreak[4]
\[\boldsymbol n_2=(-1,1,0).\]

\Needspace{5\baselineskip}
设二面角 $B-PC-A$ 的大小为 $\varphi$，则
\nopagebreak[4]
\[\cos\varphi=\frac{|\boldsymbol n_1\cdot\boldsymbol n_2|}{|\boldsymbol n_1||\boldsymbol n_2|}
=\frac{|{-1}|}{\sqrt5\cdot\sqrt2}
=\frac1{\sqrt{10}}.\]
\Needspace{5\baselineskip}
于是
\nopagebreak[4]
\[\sin\varphi=\sqrt{1-\frac1{10}}=\frac3{\sqrt{10}},\]
\Needspace{5\baselineskip}
所以
\nopagebreak[4]
\[\tan\varphi=\frac{\sin\varphi}{\cos\varphi}=3.\]

二面角 $B-PC-A$ 的正切值为 $3$。

\gkanswer{答案（1）：见解析版}

\gkanswer{答案（2）：$3$。由 $PC\perp BD$ 定出 $AB=2$，两半平面法向量可取 $(1,0,2)$、$(-1,1,0)$，夹角余弦 $\dfrac1{\sqrt{10}}$，正切为 $3$。}

\end{answer}
\begin{answer}
\gksolhead{19}{}

\textbf{题目}：

设数列 $\{a_n\}$ 的前 $n$ 项和为 $S_n$，满足 $2S_n=a_{n+1}-2^{n+1}+1$，$n\in\mathbf N^*$，且 $a_1$，$a_2+5$，$a_3$ 成等差数列。

（1） 求 $a_1$ 的值；

（2） 求数列 $\{a_n\}$ 的通项公式；

（3） 证明：对一切正整数 $n$，有 $\dfrac1{a_1}+\dfrac1{a_2}+\dfrac1{a_3}+\cdots+\dfrac1{a_n}<\dfrac32$。

\Needspace{4\baselineskip}
\textbf{解析}

\textbf{（1）}

\Needspace{5\baselineskip}
在 $2S_n=a_{n+1}-2^{n+1}+1$ 中令 $n=1$。因 $S_1=a_1$，
\nopagebreak[4]
\[2a_1=a_2-2^2+1=a_2-3,\]
\Needspace{5\baselineskip}
故
\nopagebreak[4]
\[a_2=2a_1+3.\]

\Needspace{5\baselineskip}
再令 $n=2$。因 $S_2=a_1+a_2$，
\nopagebreak[4]
\[2(a_1+a_2)=a_3-2^3+1=a_3-7,\]
\Needspace{5\baselineskip}
故
\nopagebreak[4]
\[a_3=2a_1+2a_2+7.\]

\Needspace{5\baselineskip}
由 $a_1$，$a_2+5$，$a_3$ 成等差数列，
\nopagebreak[4]
\[2(a_2+5)=a_1+a_3.\]
\Needspace{5\baselineskip}
把 $a_2$、$a_3$ 的表达式代入：
\nopagebreak[4]
\[2a_2+10=a_1+(2a_1+2a_2+7)=3a_1+2a_2+7.\]
\Needspace{5\baselineskip}
两边消去 $2a_2$：
\nopagebreak[4]
\[10=3a_1+7,\qquad 3a_1=3,\]
\Needspace{5\baselineskip}
所以
\nopagebreak[4]
\[a_1=1.\]

\textbf{（2）}

\Needspace{5\baselineskip}
当 $n\geqslant2$ 时，把已知式中的 $n$ 换成 $n-1$：
\nopagebreak[4]
\[2S_{n-1}=a_n-2^n+1.\]
\Needspace{5\baselineskip}
与原式
\nopagebreak[4]
\[2S_n=a_{n+1}-2^{n+1}+1\]
\Needspace{5\baselineskip}
相减。左边为 $2(S_n-S_{n-1})=2a_n$，故
\nopagebreak[4]
\[2a_n=a_{n+1}-a_n-2^{n+1}+2^n,\]
\Needspace{5\baselineskip}
整理得
\nopagebreak[4]
\[a_{n+1}=3a_n+2^n,\qquad n\geqslant2.\]
\Needspace{5\baselineskip}
又由（1），$a_1=1$，$a_2=2a_1+3=5$，而
\nopagebreak[4]
\[3a_1+2^1=3+2=5,\]
\Needspace{5\baselineskip}
故递推式
\nopagebreak[4]
\[a_{n+1}=3a_n+2^n\]
对 $n\geqslant1$ 都成立。

\Needspace{5\baselineskip}
两边同除以 $3^{n+1}$：
\nopagebreak[4]
\[\frac{a_{n+1}}{3^{n+1}}=\frac{a_n}{3^n}+\frac13\left(\frac23\right)^n.\]
\Needspace{5\baselineskip}
令 $b_n=\dfrac{a_n}{3^n}$，则
\nopagebreak[4]
\[b_{n+1}-b_n=\frac13\left(\frac23\right)^n.\]
\Needspace{5\baselineskip}
累加，$b_1=\dfrac{a_1}{3}=\dfrac13$，得
\nopagebreak[4]
\[\begin{aligned}
b_n&=b_1+\frac13\sum_{k=1}^{n-1}\left(\frac23\right)^k\\
&=\frac13+\frac13\cdot\frac{\frac23\left[1-\left(\frac23\right)^{n-1}\right]}{1-\frac23}\\
&=\frac13+\frac23\left[1-\left(\frac23\right)^{n-1}\right]\\
&=1-\left(\frac23\right)^n.
\end{aligned}\]
\Needspace{5\baselineskip}
所以
\nopagebreak[4]
\[a_n=3^n b_n=3^n-2^n.\]
检验：$a_1=3-2=1$，$a_2=9-4=5$，$a_3=27-8=19$，且 $1,10,19$ 成等差数列，符合。

\textbf{（3）}

\Needspace{5\baselineskip}
由（2），$a_k=3^k-2^k$。对 $k\geqslant2$，
\nopagebreak[4]
\[3^k>2^{k+1},\]
\Needspace{5\baselineskip}
因为
\nopagebreak[4]
\[\left(\frac32\right)^k\geqslant\left(\frac32\right)^2=\frac94>2.\]
\Needspace{5\baselineskip}
于是
\nopagebreak[4]
\[3^k-2^k>2^{k+1}-2^k=2^k,\]
\Needspace{5\baselineskip}
故
\nopagebreak[4]
\[\frac1{a_k}=\frac1{3^k-2^k}<\frac1{2^k}.\]

\Needspace{5\baselineskip}
因此
\nopagebreak[4]
\[\begin{aligned}
\sum_{k=1}^n\frac1{a_k}
&=\frac1{a_1}+\sum_{k=2}^n\frac1{a_k}\\
&<1+\sum_{k=2}^n\frac1{2^k}\\
&<1+\sum_{k=2}^{\infty}\frac1{2^k}\\
&=1+\frac12=\frac32.
\end{aligned}\]
\Needspace{5\baselineskip}
其中 $k=1$ 时 $\dfrac1{a_1}=1$ 单独保留；第二个不等号来自等比级数
\nopagebreak[4]
\[\sum_{k=2}^{\infty}\frac1{2^k}=\frac{\frac14}{1-\frac12}=\frac12.\]
\Needspace{5\baselineskip}
故对一切正整数 $n$，
\nopagebreak[4]
\[\frac1{a_1}+\frac1{a_2}+\cdots+\frac1{a_n}<\frac32.\]

\gkanswer{答案（1）：$a_1=1$。由 $n=1,2$ 得 $a_2=2a_1+3$、$a_3=2a_1+2a_2+7$，代入等差条件即得。}

\gkanswer{答案（2）：$a_n=3^n-2^n$。作差得 $a_{n+1}=3a_n+2^n$，化归等比数列求解。}

\gkanswer{答案（3）：见解析版}

\end{answer}
\begin{answer}
\gksolhead{20}{}

\textbf{题目}：

在平面直角坐标系 $xOy$ 中，已知椭圆 $C:\dfrac{x^2}{a^2}+\dfrac{y^2}{b^2}=1$（$a>b>0$）的离心率 $e=\sqrt{\dfrac23}$，且椭圆 $C$ 上的点到点 $Q(0,2)$ 的距离的最大值为 $3$。

（1） 求椭圆 $C$ 的方程；

（2） 在椭圆 $C$ 上，是否存在点 $M(m,n)$，使得直线 $l:mx+ny=1$ 与圆 $O:x^2+y^2=1$ 相交于不同的两点 $A$，$B$，且 $\triangle OAB$ 的面积最大？若存在，求出点 $M$ 的坐标及相对应的 $\triangle OAB$ 的面积；若不存在，请说明理由。

\Needspace{4\baselineskip}
\textbf{解析}

\textbf{（1）}

\Needspace{5\baselineskip}
设椭圆的半焦距为 $c$，则
\nopagebreak[4]
\[e^2=\frac{c^2}{a^2}=\frac23,\qquad c^2=\frac23a^2.\]
\Needspace{5\baselineskip}
于是
\nopagebreak[4]
\[b^2=a^2-c^2=a^2-\frac23a^2=\frac13a^2.\]

\Needspace{5\baselineskip}
设 $P(x,y)$ 是椭圆上一点，则
\nopagebreak[4]
\[\frac{x^2}{a^2}+\frac{y^2}{b^2}=1,\qquad
x^2=a^2\left(1-\frac{y^2}{b^2}\right).\]
\Needspace{5\baselineskip}
因 $b^2=\dfrac{a^2}3$，得
\nopagebreak[4]
\[x^2=a^2-3y^2.\]
\Needspace{5\baselineskip}
点 $P$ 到 $Q(0,2)$ 的距离的平方为
\nopagebreak[4]
\[\begin{aligned}
d^2&=x^2+(y-2)^2\\
&=a^2-3y^2+y^2-4y+4\\
&=a^2+4-4y-2y^2\\
&=a^2+6-2(y+1)^2.
\end{aligned}\]
其中 $y\in[-b,b]$。

\Needspace{5\baselineskip}
若 $b<1$，则 $-1\notin[-b,b]$，$d^2$ 在 $[-b,b]$ 上递减，最大值在 $y=-b$ 处：
\nopagebreak[4]
\[d^2_{\max}=a^2+4+4b-2b^2=3b^2+4+4b-2b^2=(b+2)^2<9,\]
与最大值为 $3$（即 $d^2_{\max}=9$）矛盾。

\Needspace{5\baselineskip}
故 $b\geqslant1$，此时 $-1\in[-b,b]$，$d^2$ 在 $y=-1$ 处取得最大值：
\nopagebreak[4]
\[d^2_{\max}=a^2+6-0=a^2+6.\]
\Needspace{5\baselineskip}
由 $d_{\max}=3$ 得
\nopagebreak[4]
\[a^2+6=9,\qquad a^2=3,\]
\Needspace{5\baselineskip}
从而
\nopagebreak[4]
\[b^2=\frac13a^2=1.\]
\Needspace{5\baselineskip}
椭圆方程为
\nopagebreak[4]
\[\frac{x^2}3+y^2=1.\]
此时 $b=1$，$y=-1$ 是区间端点，点 $(0,-1)$ 到 $Q(0,2)$ 的距离为 $3$，确能取到。

\textbf{（2）}

\Needspace{5\baselineskip}
设圆心 $O$ 到直线 $l:mx+ny=1$ 的距离为 $d$。由点到直线的距离公式，
\nopagebreak[4]
\[d=\frac1{\sqrt{m^2+n^2}}.\]
\Needspace{5\baselineskip}
直线与单位圆相交于不同两点，需 $d<1$；此时弦长
\nopagebreak[4]
\[|AB|=2\sqrt{1-d^2}.\]
\Needspace{5\baselineskip}
$\triangle OAB$ 的面积为
\nopagebreak[4]
\[S=\frac12|AB|\cdot d=d\sqrt{1-d^2}.\]
\Needspace{5\baselineskip}
于是
\nopagebreak[4]
\[S^2=d^2(1-d^2)\leqslant\left(\frac{d^2+1-d^2}{2}\right)^2=\frac14,\]
\Needspace{5\baselineskip}
等号当且仅当
\nopagebreak[4]
\[d^2=1-d^2,\qquad d^2=\frac12.\]
\Needspace{5\baselineskip}
所以 $S$ 的最大值为 $\dfrac12$，且取最大值时
\nopagebreak[4]
\[m^2+n^2=\frac1{d^2}=2.\]

\Needspace{5\baselineskip}
又 $M(m,n)$ 在椭圆 $C$ 上，
\nopagebreak[4]
\[\frac{m^2}3+n^2=1,\]
\Needspace{5\baselineskip}
即
\nopagebreak[4]
\[m^2+3n^2=3.\]
\Needspace{5\baselineskip}
与 $m^2+n^2=2$ 联立，相减得
\nopagebreak[4]
\[2n^2=1,\qquad n^2=\frac12,\]
\Needspace{5\baselineskip}
从而
\nopagebreak[4]
\[m^2=2-n^2=\frac32.\]
\Needspace{5\baselineskip}
所以
\nopagebreak[4]
\[m=\pm\frac{\sqrt6}2,\qquad n=\pm\frac{\sqrt2}2,\]
\Needspace{5\baselineskip}
四个点
\nopagebreak[4]
\[\left(\frac{\sqrt6}2,\frac{\sqrt2}2\right),\quad
\left(\frac{\sqrt6}2,-\frac{\sqrt2}2\right),\quad
\left(-\frac{\sqrt6}2,\frac{\sqrt2}2\right),\quad
\left(-\frac{\sqrt6}2,-\frac{\sqrt2}2\right)\]
都满足条件。

\Needspace{5\baselineskip}
检验：$m^2+n^2=\dfrac32+\dfrac12=2$，故 $d=\dfrac1{\sqrt2}<1$，直线与圆确有两个不同交点；此时
\nopagebreak[4]
\[|AB|=2\sqrt{1-\frac12}=\sqrt2,\qquad
S=\frac12\cdot\sqrt2\cdot\frac1{\sqrt2}=\frac12.\]

\Needspace{5\baselineskip}
故存在这样的点 $M$，其坐标为
\nopagebreak[4]
\[\left(\pm\frac{\sqrt6}2,\pm\frac{\sqrt2}2\right),\]
相应的 $\triangle OAB$ 面积的最大值为 $\dfrac12$。

\gkanswer{答案（1）：$\dfrac{x^2}3+y^2=1$。$e^2=\dfrac23$ 给 $b^2=\dfrac{a^2}3$，距离平方的最大值为 $a^2+6=9$。}

\gkanswer{答案（2）：存在，$M$ 为 $\left(\pm\dfrac{\sqrt6}2,\pm\dfrac{\sqrt2}2\right)$（四个点），$\triangle OAB$ 面积的最大值为 $\dfrac12$。面积取最大值时圆心到直线距离为 $\dfrac1{\sqrt2}$，即 $m^2+n^2=2$，再与椭圆方程联立。}

\end{answer}
\begin{answer}
\gksolhead{21}{}

\textbf{题目}：

设 $a<1$，集合 $A=\{x\in\mathbf R\mid x>0\}$，$B=\{x\in\mathbf R\mid 2x^2-3(1+a)x+6a>0\}$，$D=A\cap B$。

（1） 求集合 $D$（用区间表示）；

（2） 求函数 $f(x)=2x^3-3(1+a)x^2+6ax$ 在 $D$ 内的极值点。

\Needspace{4\baselineskip}
\textbf{解析}

\textbf{（1）}

\Needspace{5\baselineskip}
记
\nopagebreak[4]
\[q(x)=2x^2-3(1+a)x+6a.\]
\Needspace{5\baselineskip}
其判别式为
\nopagebreak[4]
\[\begin{aligned}
\Delta&=9(1+a)^2-48a\\
&=9a^2+18a+9-48a\\
&=9a^2-30a+9\\
&=3(3a-1)(a-3).
\end{aligned}\]
\Needspace{5\baselineskip}
两根为
\nopagebreak[4]
\[x_{1,2}=\frac{3(1+a)\mp\sqrt{\Delta}}4,\qquad x_1<x_2.\]
\Needspace{5\baselineskip}
因二次项系数 $2>0$，当 $\Delta>0$ 时
\nopagebreak[4]
\[B=(-\infty,x_1)\cup(x_2,+\infty);\]
当 $\Delta=0$ 时 $B=\mathbf R\setminus\{x_1\}$；当 $\Delta<0$ 时 $B=\mathbf R$。

由 $a<1$ 分情形讨论。

\Needspace{5\baselineskip}
当 $a<0$ 时，$\Delta>0$，且两根之积 $3a<0$，故
\nopagebreak[4]
\[x_1<0<x_2.\]
\Needspace{5\baselineskip}
于是
\nopagebreak[4]
\[D=A\cap B=(0,+\infty)\cap\left((-\infty,x_1)\cup(x_2,+\infty)\right)=(x_2,+\infty).\]

\Needspace{5\baselineskip}
当 $a=0$ 时，
\nopagebreak[4]
\[q(x)=2x^2-3x=x(2x-3),\]
\Needspace{5\baselineskip}
$q(x)>0$ 的解为 $x<0$ 或 $x>\dfrac32$，故
\nopagebreak[4]
\[D=\left(\frac32,+\infty\right).\]

\Needspace{5\baselineskip}
当 $0<a<\dfrac13$ 时，$\Delta>0$，两根之和 $\dfrac{3(1+a)}2>0$、两根之积 $3a>0$，故
\nopagebreak[4]
\[0<x_1<x_2.\]
\Needspace{5\baselineskip}
于是
\nopagebreak[4]
\[D=(0,x_1)\cup(x_2,+\infty).\]

\Needspace{5\baselineskip}
当 $a=\dfrac13$ 时，$\Delta=0$，
\nopagebreak[4]
\[q(x)=2x^2-4x+2=2(x-1)^2,\]
\Needspace{5\baselineskip}
$q(x)>0$ 对 $x\neq1$ 成立，故
\nopagebreak[4]
\[D=(0,1)\cup(1,+\infty).\]

\Needspace{5\baselineskip}
当 $\dfrac13<a<1$ 时，$\Delta<0$，$q(x)>0$ 对一切实数 $x$ 成立，故
\nopagebreak[4]
\[D=(0,+\infty).\]

\Needspace{5\baselineskip}
综上：
\nopagebreak[4]
\[D=\begin{cases}
(x_2,+\infty), & a<0,\\[1mm]
\left(\dfrac32,+\infty\right), & a=0,\\[2mm]
(0,x_1)\cup(x_2,+\infty), & 0<a<\dfrac13,\\[1mm]
(0,1)\cup(1,+\infty), & a=\dfrac13,\\[1mm]
(0,+\infty), & \dfrac13<a<1,
\end{cases}\]
\Needspace{5\baselineskip}
其中
\nopagebreak[4]
\[x_{1,2}=\frac{3(1+a)\mp\sqrt{9a^2-30a+9}}4.\]

\textbf{（2）}

\Needspace{5\baselineskip}
求导：
\nopagebreak[4]
\[\begin{aligned}
f'(x)&=6x^2-6(1+a)x+6a\\
&=6\left[x^2-(1+a)x+a\right]\\
&=6(x-1)(x-a).
\end{aligned}\]
临界点为 $x=a$ 与 $x=1$。按（1）的各个情形判断它们是否落在 $D$ 内，并考察 $f'$ 的符号变化。

\Needspace{5\baselineskip}
当 $a<0$ 时，$D=(x_2,+\infty)$。由
\nopagebreak[4]
\[q(1)=2-3(1+a)+6a=3a-1<0,\]
\Needspace{5\baselineskip}
且 $q(x)>0$ 在 $(x_2,+\infty)$ 上成立，可知 $1<x_2$。于是 $D$ 内每一点都满足 $x>1>a$，从而
\nopagebreak[4]
\[f'(x)=6(x-1)(x-a)>0,\]
$f$ 在 $D$ 上单调递增，无极值点。

\Needspace{5\baselineskip}
当 $a=0$ 时，$D=\left(\dfrac32,+\infty\right)$，在 $D$ 内 $x>1>0$，
\nopagebreak[4]
\[f'(x)=6x(x-1)>0,\]
$f$ 单调递增，无极值点。

\Needspace{5\baselineskip}
当 $0<a<\dfrac13$ 时，$D=(0,x_1)\cup(x_2,+\infty)$，且
\nopagebreak[4]
\[0<a<x_1<1<x_2.\]
\Needspace{5\baselineskip}
这里 $a<x_1$ 是因为
\nopagebreak[4]
\[q(a)=2a^2-3(1+a)a+6a=a(3-a)>0,\]
故 $a\in B$，而 $a<1<x_2$，只能 $a\in(0,x_1)$。

在 $(0,a)$ 上，$x-1<0$、$x-a<0$，故 $f'>0$；在 $(a,x_1)$ 上，$x-1<0$、$x-a>0$，故 $f'<0$。所以 $x=a$ 是极大值点。

在 $(x_2,+\infty)$ 上，$x>1>a$，$f'>0$，没有极值点。

\Needspace{5\baselineskip}
故此时只有极大值点
\nopagebreak[4]
\[x=a.\]

\Needspace{5\baselineskip}
当 $a=\dfrac13$ 时，$D=(0,1)\cup(1,+\infty)$，
\nopagebreak[4]
\[f'(x)=6(x-1)\left(x-\frac13\right).\]
\Needspace{5\baselineskip}
在 $\left(0,\dfrac13\right)$ 上 $f'>0$，在 $\left(\dfrac13,1\right)$ 上 $f'<0$，故 $x=\dfrac13$ 是极大值点；$x=1\notin D$，不是 $D$ 内的极值点。故此时只有极大值点
\nopagebreak[4]
\[x=\frac13.\]

\Needspace{5\baselineskip}
当 $\dfrac13<a<1$ 时，$D=(0,+\infty)$。在 $(0,a)$ 上 $f'>0$，在 $(a,1)$ 上 $f'<0$，在 $(1,+\infty)$ 上 $f'>0$。故
\nopagebreak[4]
\[x=a\ \text{\text{是极大值点}},\qquad x=1\ \text{\text{是极小值点}}.\]

\Needspace{5\baselineskip}
综上：
\nopagebreak[4]
\[a\leqslant0\ \text{\text{时无极值点}};\qquad
0<a<\frac13\ \text{\text{时极大值点为} }x=a;\]
\nopagebreak[4]
\[a=\frac13\ \text{\text{时极大值点为} }x=\frac13;\qquad
\frac13<a<1\ \text{\text{时极大值点为} }x=a,\ \text{\text{极小值点为} }x=1.\]

\gkanswer{答案（1）：记 $\Delta=9a^2-30a+9=3(3a-1)(a-3)$，$x_{1,2}=\dfrac{3(1+a)\mp\sqrt\Delta}4$（$x_1<x_2$）。当 $a<0$ 时 $D=(x_2,+\infty)$；当 $a=0$ 时 $D=\left(\dfrac32,+\infty\right)$；当 $0<a<\dfrac13$ 时 $D=(0,x_1)\cup(x_2,+\infty)$；当 $a=\dfrac13$ 时 $D=(0,1)\cup(1,+\infty)$；当 $\dfrac13<a<1$ 时 $D=(0,+\infty)$。}

\gkanswer{答案（2）：$a\leqslant0$ 时无极值点；$0<a<\dfrac13$ 时极大值点为 $x=a$；$a=\dfrac13$ 时极大值点为 $x=\dfrac13$；$\dfrac13<a<1$ 时极大值点为 $x=a$、极小值点为 $x=1$。依据 $f'(x)=6(x-1)(x-a)$ 及 $a$、$1$ 是否落在 $D$ 内判定。}

\end{answer}